% Määrittele avainsanat main.tex:ssä ja poista seuraava rivi
%\avainsanat{}
\begin{thesisabstract}{finnish}
\tiivistelmametadata
% Voit myös laittaa tähän version tiivistelmästä, joka sisältää kappalejakoja
% ja erikoismerkkejä. Sisällön tulisi kuitenkin olla muutoin identtinen
% metadatassa olevan version kanssa.

Oulun yliopiston tietojärjestelmäkurssin opiskelijat kohtaavat usein vaikeuksia työkalujensa asennuksessa ja joutuvat turvautumaan yksinkertaisiin tekstilähtöihin etsiessään virheitä koodistaan. Tässä opinnäytetyössä esitellään yhtenäisen, konttiteknologiaan perustuvan virheenkorjausalustan suunnittelu ja toteutus Raspberry Pi RP2040- ja RP2350-mikro-ohjaimille. Alusta hyödyntää Docker-konttitekniikkaa, virtuaalikoneita ja verkkopohjaista IDE:tä (code-server) tarjotakseen skaalautuvan ratkaisun, joka takaa toistettavan opetusympäristön, jossa
ympäristövirheet eliminoidaan.

Tämän työn keskeinen saavutus on Pico-Debug-Dashboardin kehittäminen, joka on räätälöity Visual Studio Code -laajennus, joka yhdistää CMaken, OpenOCD:n ja GDB:n monimutkaisuuden keskitettyyn käyttöliittymään. Alusta tukee myös useita virheenkorjausmenetelmiä mukaan lukien fyysiset laitteistokoettimet ja virtuaaliset ohjelmistokoettimet (pico-debug), sekä räätälöidyn integroidun kotelosuunnittelun, joka sopii alustan modulaarisiin tarpeisiin. Jotta voitiin arvioida alustan tehokkuutta opetuskäytössä Oulun yliopiston tietojärjestelmäkurssilla tehtiin vertaileva analyysi konttitekniikkaan ja virtuaalikoneisiin perustuvien toteutusten välillä.

Arviointi suoritettiin käyttämällä sekä määrällisiä mittareita, kuten käyttöönoton viivettä ja flash-muistin siirtonopeutta, että laadullisia mittareita, joissa keskityttiin aloittelijaystävällisyyteen ja ergonomiaan. Tulokset osoittavat, että konttitekniikka lyhentää merkittävästi ympäristön kylmäkäynnistysaikoja ja tarjoaa erinomaisen laitteistoyhteyden vakauden. Arvioiduista virheenkorjausantureista Second Pico -laitteistototeutus todettiin sopivimmaksi laajamittaiseen käyttöönottoon luokkahuoneissa, koska se tarjoaa tasapainoisen tuen, edullisen hinnan ja kaksisydämisen virheenkorjausominaisuuden.
\end{thesisabstract}
